\documentclass[11pt]{article}
\usepackage[a4paper,margin=2.2cm]{geometry}
\usepackage{xcolor}
\usepackage{fontspec}
\usepackage{xeCJK}
\usepackage{enumitem}
\setmainfont{Times New Roman}
\setCJKmainfont{SimSun}
\setlist[itemize]{leftmargin=1.4em,itemsep=0.35em,topsep=0.35em}
\setlength{\parindent}{0pt}
\setlength{\parskip}{0.55em}
\pagestyle{plain}
\begin{document}
{\color[HTML]{1F5FD2}\LARGE\bfseries 社交技能训练课程设计\par}
{\color[HTML]{667085}\small 星轨原创教育资源：用于家庭与学校共同制定支持计划。\par}
\vspace{0.8em}
{\color[HTML]{111827}\large\bfseries 课程目标\par}
\begin{itemize}
\item 社交技能训练要把抽象能力转化为可观察行为，例如发起互动、等待、回应、协商、拒绝和修复冲突。
\item 目标应选自儿童真实社交场景，避免只在训练室中完成而无法迁移。
\end{itemize}
{\color[HTML]{111827}\large\bfseries 教学流程\par}
\begin{itemize}
\item 每次只教一个核心技能，可用图片或视频解释，再示范正确与不合适的例子，随后进行角色扮演。
\item 练习后尽快安排真实活动，如课间、午餐或合作游戏，让儿童在自然情境中使用同一技能。
\end{itemize}
{\color[HTML]{111827}\large\bfseries 泛化应用\par}
\begin{itemize}
\item 学校和家庭应使用一致的提示词，帮助儿童在不同地点识别同一个社交要求。
\item 如果儿童出现回避，先降低社交复杂度，再逐步增加伙伴数量、等待时间或对话轮次。
\end{itemize}
\vfill
{\color[HTML]{667085}\footnotesize 本材料为应用内原创教育内容，不能替代医疗、法律或正式评估建议。\par}
\end{document}
